\section{Experimental Setup and Metrics}

\subsection{Compute Environment and Memory Limits}
\textbf{Configuration~A (GPU)} ran on a free Kaggle notebook with one NVIDIA Tesla T4 (15\,GB; the notebook exposes two, only one is used), PyTorch~2.11 (CUDA~13 build), a recent Hugging Face Transformers~5.x~\cite{wolf2020transformers} and PEFT~0.21 for training and for reference checks. Training used fp16 autocast with fp32 master weights and LoRA parameters; \emph{all inference, evaluation and benchmarks run in fp32} so that exact-equality tests are meaningful. The 1.7B target needs about 6.85\,GB for fp32 weights, which together with activations left too little room for a micro-batch of 4 on the first attempt (an out-of-memory error), hence micro-batch 1 with 16 accumulation steps.

\textbf{Configuration~B (CPU)} ran on a laptop: Intel Core i5-1135G7 (4 cores / 8 threads), 16\,GB RAM, Windows~11, no GPU, PyTorch~2.14 (CPU build), Transformers~5.22 and PEFT~0.21. Everything uses fp32 on four threads, because CPU bf16 matrix multiplication is slower than fp32 on this processor and would add rounding noise to the exact-equality tests. The 360M target needs roughly 1.45\,GB for weights and, during LoRA training with activations and AdamW state, a peak working set of about 4\,GB; the 135M draft needs about 0.55\,GB. Training time on this machine is the practical limit that fixed the model size (Section~III-C).

Both configurations use exactly the same data split, prompts, parser, metrics and code; only the model pair, device and training schedule differ. The CPU configuration also serves as the model of the locally hosted application.

\subsection{Classification Metrics}
We report accuracy, macro-F1 and per-class precision, recall and F1 with support, plus confusion matrices, for all test criteria and separately for inclusion and exclusion criteria. Macro-F1 averages over the classes that occur either in the gold labels or in the predictions of the scope under consideration; a class absent from both is skipped, so the missing \emph{excluded} class in the test split does not silently count as a zero. An \emph{invalid} model output is counted as wrong and is never mapped onto a label. Because the test split contains only 8 patients whose rows are correlated, we also give 95\% patient-level bootstrap intervals (2\,000 resamples of patients).

\subsection{Evidence-Sentence Metrics}
Evidence is scored on sentence IDs as set overlap. \emph{Micro} scores pool true/false positives and false negatives over all records. \emph{Macro} scores average per-record precision, recall and F1 in two ways: (a) over records with at least one gold sentence (``evidence-bearing''), where predicting nothing yields precision 0 by convention, and (b) over all records, where an empty gold set matched by an empty prediction counts as a perfect record. Records with empty gold evidence are additionally reported on their own: the fraction for which the model also cited nothing and the number of spurious citations it made. About \NGoldEmpty{} of the \NGoldNon{}+\NGoldEmpty{} test records have no annotated evidence, so the policy matters and we show both views. Absence of cited evidence is never converted into an eligibility judgement.

\subsection{Decoding Settings}
Every accuracy number uses deterministic greedy decoding (argmax), at most 48 new tokens, stopping at the end-of-turn token, with the same prompt for the base and the adapted model. The base model is evaluated zero-shot with the identical prompt and parser on exactly the same validation and test rows; it receives no few-shot examples and no adapter. Hyper-parameters were chosen once from the configuration in Section~III-D and not tuned on the test split.

\subsection{Inference Benchmarks}
\emph{KV cache.} For prompt lengths $\{64,256,512\}$ tokens (real prompts built from test criteria, concatenated and cut) and output lengths $\{16,64,128\}$, the uncached reference, the cached path and the cached path with the tiled kernel generate with EOS suppressed so that all paths emit exactly the same number of tokens. These benchmarks were run on the GPU of Configuration~A (1.7B target); the laptop CPU was used for training, evaluation and the live comparison inside the application (Section~\ref{sec:appres}), not for stored benchmark tables. Each benchmark point runs in a fresh process with one warm-up generation and three timed repetitions; we report medians of prefill time (time to the first token), decode time, tokens per second and peak memory growth. On the GPU the memory figure is the exact CUDA-allocator peak minus the allocation before the call (\texttt{torch.cuda.max\_memory\_allocated}); in the application on a CPU it is the resident-set size sampled every 2\,ms while generating (peak minus the value before the call), a proxy because the CPU allocator has no peak counter. Greedy token sequences must be identical across paths.

\emph{Attention.} Shapes follow the 360M model (15 query heads, 5 KV heads, $d=64$) for $T\in\{128,\dots,4096\}$ on the GPU, in prefill mode ($T_q=T_k=T$) and decode mode ($T_q=1$ against $T$ cached keys), with five repetitions per point and the same fresh-process memory probe. Errors are the maximum and mean absolute difference to the reference output.

\emph{Speculative decoding.} Thirty-two test-criterion prompts (sampled with a fixed seed; only speed and token equality are measured, no model selection uses them) are decoded by cached greedy target decoding and by our speculative decoder for $k\in\{2,4,6,8\}$, once with the LoRA-adapted draft and once with the un-adapted draft. We report token equality, acceptance rate (accepted drafted tokens / drafted tokens), mean accepted tokens per round, target passes per emitted token, mean latency, tokens per second and speed-up. Two output regimes are used: the \emph{natural} task output (about 15 tokens of JSON) and a \emph{forced} 64-token output (EOS suppressed) that exposes behaviour on longer generations. Hugging Face assisted generation serves as an independent equality check on a subset.
